\chapter{Synthesis strategies and top-down methods}
\label{ch:topdown}

\paragraph{In this chapter.} The two families of synthesis and the top-down methods
(ball milling, lithography/selective etching). Each method is summarised on the same card so it can
be compared later in \cref{ch:synth-compare}.

\section{Two strategies}
\begin{center}
\begin{tikzpicture}[node distance=6mm and 10mm, every node/.style={font=\small}]
\node[draw, rounded corners, fill=conc!10] (bulk) {Bulk};
\node[draw, rounded corners, fill=conc!10, right=of bulk] (pow) {Powder (\si{\micro m})};
\node[draw, rounded corners, fill=keyc!10, right=of pow] (np) {\textbf{Nanoparticles}};
\node[draw, rounded corners, fill=proc!10, right=of np] (cl) {Clusters};
\node[draw, rounded corners, fill=proc!10, right=of cl] (at) {Atoms};
\draw[-{Stealth}, thick, conc] (bulk) -- (pow);
\draw[-{Stealth}, thick, conc] (pow) -- (np);
\draw[-{Stealth}, thick, proc] (at) -- (cl);
\draw[-{Stealth}, thick, proc] (cl) -- (np);
\node[conc, below=2mm of pow] {top-down (breakdown)};
\node[proc, below=2mm of cl] {bottom-up (build-up)};
\end{tikzpicture}
\end{center}

\begin{keybox}[: top-down vs bottom-up \slide{29}]
\begin{tabularx}{\linewidth}{@{}lX@{}}
\textbf{Top-down} & Mechanical/chemical breakdown of bulk. \emph{Simple, but poor size control.} \\
\textbf{Bottom-up} & Atoms/ions/molecules assemble (vaporisation or dissolution $\to$ metal ions/clusters $\to$
reduction, rapid cooling, dispersion/protection agents). \emph{Complex, but good size and shape control.}
\end{tabularx}
\end{keybox}

Bottom-up methods are split in this course into \textbf{gas phase} (\cref{ch:gas}) and \textbf{liquid phase}
(\cref{ch:liquid}).

\section{Ball milling}
\begin{cardbox}{Ball milling \slide{31--32}}
\textbf{Principle.} Powder + hard balls in a rotating (planetary) jar. Repeated impacts cause elastic $\to$ plastic
$\to$ shear deformation, then fracture, amorphisation and even chemical reactions (mechanochemistry).\\
\textbf{Controls size/shape:} milling time, energy (speed), ball/powder ratio. Size control is poor.\\
\textbf{Typical products:} oxide powders, mixed oxides (e.g.\ CuO--CeO$_2$, CeO$_2$--ZrO$_2$ shown in
STEM/EDX maps on slide 32).
\end{cardbox}
\proscons{
\begin{itemize}
  \item Cheap, fast, easy.
  \item \textbf{No solvents} (green).
  \item Scalable.
\end{itemize}}{
\begin{itemize}
  \item Poor size control (broad distribution).
  \item Possible surface contamination (from balls/jar).
  \item Introduction of defects (dislocations).
\end{itemize}}
\textbf{Choose it when:} you need large amounts of a powder, size distribution is not critical, and defects
can even be useful (e.g.\ for catalysis).

\section{Lithography and selective etching}
\begin{cardbox}{Selective etching of Si \slide{33--36}}
\textbf{Principle.} Si wafer with SiO$_2$ mask $\to$ lithography $\to$ Si etch in TMAH (pyramidal pits)
$\to$ photo-assisted (electrochemical) etching makes deep channels $\to$ oxidation $\to$ etching in TMAH
$\to$ etching in HF. Result: regular arrays of $\sim4\times10^6$ channels/cm$^2$ ($\phi\approx\SI{3.3}{\micro m}$).\\
\textbf{Functionalisation.} Surface OH groups (after thermal treatment $T_1$) react with a metal precursor
ML$_m$ (grafting, \slide{35}); a second treatment $T_2$ gives a thin MO$_x$ layer, onto which metal nanoparticles
can be deposited. Examples: Au/TiO$_2$ and RhPd/CeO$_2$ coated on the channel walls (\emph{micro-monoliths}
for catalysis, \slide{36}).
\end{cardbox}
\proscons{
\begin{itemize}
  \item Perfect control of geometry/pattern at the \si{\micro m} scale.
  \item Very high, ordered surface area; reproducible.
\end{itemize}}{
\begin{itemize}
  \item Multi-step, expensive (clean room).
  \item Limited to substrates such as Si; the nanoparticles themselves still need a bottom-up step.
\end{itemize}}
\textbf{Choose it when:} you need an ordered, patterned support (micro-reactors, sensors, structured catalysts).
